\subsection{Conclusiones}

En este trabajo se ha analizado el juego de las monedas, donde dos jugadores (Sophia y Mateo) eligen monedas de una fila
según ciertas reglas y Sophia, aplicando un enfoque greedy, busca ganar siempre. A través de la implementación del algoritmo 
propuesto, hemos logrado simular el comportamiento de los jugadores demostrando que todas las situaciones\footnote{ Excepto el
caso en que las monedas inicialmente son una cantidad par y son todas iguales. No hay forma de que esto no termine en empate.} 
desembocan  en el triunfo de Sophia. Esto lo logramos aplicando iterativamente una simple regla: dadas las monedas disponibles
en sus respectivos turnos, Sophia elige la moneda de mayor valor y Mateo la de menor valor.

La solución propuesta cumple con la función de optimizar la selección de monedas para Sophia, demostrando como los algoritmos greedy
pueden ser efectivos en este tipo de situaciones. Pero, a pesar de ser eficiente y práctico, este enfoque tiene sus limitaciones. Una
de ellas sucede si Sophia busca maximizar su suma acumulada y ganar. Con el algoritmo planteado dicho problema no se puede resolver
de forma óptima y aunque nos sea costoso, podríamos usar backtracking en este caso.